\section{从热词到产业拐点：为什么 2025 年成了 Agent 年}

\subsection{热度背后的三条现实支撑}
主持人在开场时给出的判断很直接：2025 年并不是 Agent 概念第一次出现，但它第一次同时满足了资本、产品和研究三个方向的共振。市场规模翻倍、头部公司推进 IPO、学术论文与开源系统同步爆发，这些信号使得 Agent 从 ``demo 概念'' 变成了企业和研究机构共同下注的赛道。

\begin{importantbox}{为什么这次不是短期噪声}
这一轮 Agent 热潮的不同之处，在于它不再只依赖一个模型榜单或一次产品发布，而是出现了完整链条：
\begin{itemize}
  \item 上游有更强的 LLM、reasoning model 与 computer use 模型；
  \item 中游有 Cursor、Claude Code、mobile agent、deep analysis 这类实际系统；
  \item 下游有企业数据分析、软件工程、GUI 自动化等明确使用场景。
\end{itemize}
\end{importantbox}

与前几轮 ``prompt engineering 热'' 相比，这场 panel 更关心的问题不是 ``怎么写 prompt''，而是 ``怎么把模型变成能长期工作、能和环境交互、能承接工作流的系统能力''。这也是 ``agent'' 与 ``agentic'' 被反复区分的背景。

\begin{knowledgebox}{Agent 与 Agentic 的差别}
嘉宾们反复强调，早期 Agent 更多是在 prompting 框架上堆出工作流；而 Agentic 更强调把规划、交互、恢复、探索等能力系统化，甚至进一步内化进模型参数。换句话说，前者偏 ``外接脚手架''，后者偏 ``原生能力建设''。
\end{knowledgebox}

\subsection{讨论对象已经从单轮对话变成长期生产力}
张少磊给出的例子很有代表性。他们团队把 attention 从一问一答的聊天模式，转向了 data science 自动化场景：企业有大量数据库和分析需求，传统流程需要人工编写 SQL、清洗数据、反复验证结果，而 agentic 系统能把编码、分析、评估与洞察提炼连成闭环。讨论真正切中了一个关键变化：Agent 不再只是回答问题，而是在替人完成工作链路中的多个步骤。

\begin{warningbox}{不要把 Agent 等同于一个会调用工具的聊天机器人}
如果只把 Agent 理解成 ``LLM + tool call''，很容易低估它的工程复杂度。真正难的部分在于：任务拆解是否合理、工具空间是否稳定、记忆能否跨轮保留、错误能否恢复、环境反馈能否形成训练信号。
\end{warningbox}

\subsection{本章小结}
这场 panel 的起点非常清楚：2025 年之所以成为 Agent 年，不是因为概念更新，而是因为模型能力、系统工程和产业需求第一次形成了闭环。接下来所有讨论，都是围绕这个闭环中哪一段最短、该优先补哪一段展开的。

